% Table 5 of paper_bc_peerjcs.tex, exported for PeerJ upload by
% scripts/pubs/export_peerj_upload.py. Do not edit: edit the manuscript and re-export.
\documentclass[10pt]{article}
\usepackage[letterpaper,margin=25mm]{geometry}
\usepackage{amsmath,amssymb}
\usepackage{booktabs}
\usepackage{tabularx}
\usepackage{array}
\usepackage{multirow}
\usepackage{enumitem}
\usepackage[hidelinks]{hyperref}
\usepackage[authoryear,round]{natbib}
\begin{document}
\begin{table}[t]
\centering
\caption{Working construct dictionary: operational meanings and strongest
direct evidence for each quality property.}
\begin{tabularx}{\linewidth}{l X X}
\toprule
\textbf{Construct} & \textbf{Operational meaning} & \textbf{Strongest evidence} \\
\midrule
Fidelity / faithfulness & Whether the explanation tracks model behavior under a specified intervention or masking scheme. & Proxy tests and benchmark-grounded checks \\
Stability / robustness & Whether explanations remain consistent across reruns, perturbations, or nearby inputs. & Proxy perturbation studies and repeated-run analyses \\
Sparsity / parsimony & Whether explanations keep feature load low enough to remain cognitively manageable. & Structural summaries of explanation size or weight concentration \\
Plausibility & Whether explanations align with expert expectations or domain logic. & Expert-human review with explicit rubrics \\
Semantic alignment & Whether explanations capture meaningful conceptual or causal structure beyond structural proxy success. & Expert review or validated semantic evaluators \\
User usefulness & Whether explanations improve comprehension, decision quality, or trust calibration in an actual task. & End-user studies with task outcomes \\
\bottomrule
\end{tabularx}

\end{table}
\end{document}
